\section{Output files}

\subsection{output/ResultHydro\_pDDDDdHHhMMmSSsXXX.txt}\label{resulsubt_hydro_txt}

\subsubsection{Purpose}
This file stores the hydrodynamic results every \param{dt\_map}. To avoid any computational timestep reduction to store the results precisely every \param{dt\_map}, the results are stored when \pfg{explain algorithm}. The stored variables depend upon the chosen hyperbolic model (defined by the parameter \param{model\_hyperb}) and the \param{storeflag} parameter:
\begin{center}
\begin{tabular}{c|c|c}
\multirow{2}{*}{\param{model\_hyperb}} & \multicolumn{2}{c}{\param{storeflag} value} \\
& default & all \\
\hline
%ddp  &  &  \\
dip  &  &  \\
sp   &  &  \\
swes & $h, \overrightarrow{q}$ & $h, \overrightarrow{q}, z, \overrightarrow{u}, \mathrm{Fr}, H_s, H_t, CFL$
\end{tabular}
\pfg{complete}
\end{center}

For vectorial fields, three values are stored in each cell: the two horizontal component of the vector and the vector norm.

\begin{center}
\begin{tabular}{c|l|c}
Symbol & Variable & Type of field \\
\hline
$CFL$ & cfl & scalar \\
$\mathrm{Fr}$ & froude & scalar \\
$h$ & water\_depth & scalar \\
$H_s$ & specific\_head & scalar \\
$H_t$ & total\_head & scalar \\
$\overrightarrow{q}$ & unit\_discharge\_x, unit\_discharge\_y and	unit\_discharge\_norm & vectorial \\
$\overrightarrow{u}$ & velocity\_x,	velocity\_y	and velocity\_norm & vectorial \\
$z$ & surface\_elevation & scalar
\end{tabular}
\end{center}

\subsubsection{Format description}
The precise simulation time for storing the variables is indicated in the file name using the ISO 8601 format. The file name is \file{ResultHydro\_pDDDDdHHhMMmSSsXXX.txt} where \emph{pDDDDdHHhMMmSSsXXX} is the elapsed time from the beginning of the simulation:
\begin{itemize}
    \item DDDD is the number of days (always on 4 digits)
    \item HH is the number of hours (always on 2 digits)
    \item MM is the number of minutes (always on 2 digits)
    \item SS is the number of seconds (always on 2 digits)
    \item XXX is the number of milliseconds (always on 3 digits)
\end{itemize}

% \subsection{Example}
% \begin{verbatim}
    
% \end{verbatim}